% titel: Blatt 5, A5: Elementare Funktionen
% Teil II: der durchgerechnete Weg.

\subsection*{Gegeben und gesucht}
Zu drei Vorgaben je eine elementare Funktion mit passendem
Definitionsbereich.

\subsection*{a) Maximum bei 2, Grenzwert 1}
\emph{Gedankenschritt / Ansatz:} Eine \emph{Glocke} $9\mathrm e^{-(x-2)^2}$
hat ihr Maximum bei $x=2$, und die additive Konstante $1$ setzt das
asymptotische Niveau. Die Höhe $10$ ergibt sich aus $1+9$.
\begin{equation*}
  f(x)=1+9\,\mathrm e^{-(x-2)^2}.
\end{equation*}
Prüfung: $f(2)=1+9=10$; für $x\to\pm\infty$ ist
$\mathrm e^{-(x-2)^2}\to0$, also $f(x)\to1$. Ableitung
$f'(x)=-18(x-2)\mathrm e^{-(x-2)^2}$: Für $x<2$ ist $f'(x)>0$ (steigend),
für $x>2$ ist $f'(x)<0$ (fallend) -- ein echtes Maximum bei $x=2$.

\subsection*{b) Auf $R_{>1}$, Werte $0$ und $-1$, Pol und $-\infty$}
\emph{Gedankenschritt / Ansatz:} Gefordert $f(2)=0$, $f(3)=-1$,
$\lim_{x\searrow1}f(x)=\infty$, $\lim_{x\to\infty}f(x)=-\infty$. Unter den
einfachen elementaren Funktionen leistet $-\ln(x-1)$ den Pol bei $1$ nach
$+\infty$ und den Abfall nach $-\infty$; durch einen linearen Zusatz
$cx+d$ regulieren wir die geforderten Werte, ohne die Grenzverhaltenstypen
zu zerstören (ein linearer Term ändert die Pole/Grenzwerte nicht in ihrer
\emph{Art}).
\begin{align*}
  g(x)=-\ln(x-1)+cx+d, \qquad
  0&=g(2)=2c+d,\\
  -1&=g(3)=-\ln 2+3c+d .
\end{align*}
Aus $d=-2c$ und $-1=-\ln2+c$ folgt $c=\ln2-1$, $d=2-2\ln2$:
\begin{equation*}
  g(x)=-\ln(x-1)+(\ln2-1)\,x+2-2\ln2 .
\end{equation*}
Die vier geforderten Eigenschaften (Werte, Pol $+\infty$ bei $1$,
Grenzwert $-\infty$ am Unendlichen) sind erfüllt; der Definitionsbereich
$R_{>1}$ wird von der $-1$ im $\ln$-Argument erzwungen.

\subsection*{c) Auf $R\setminus\{-1,1\}$, fallend, Pole}
\emph{Gedankenschritt / Ansatz:} Gefordert $f(0)=0$,
$\lim_{x\to1}f(x)=\infty$, $\lim_{x\nearrow-1}f(x)=-\infty$ und streng
monoton fallend auf $R_{>1}$. Ein Nenner $x^2-1$ erzeugt die Pole bei
$\pm1$ und den Definitionsbereich; ein Zähler sollte $f(0)=0$ liefern und
das Vorzeichen über der rechten Hälfte konstant negativ machen. Kandidat
$f(x)=2x/(x^2-1)$:
\begin{equation*}
  f'(x)=\frac{2(x^2-1)-2x(2x)}{(x^2-1)^2}
  =\frac{-2x^2-2}{(x^2-1)^2}<0 \quad\text{auf }R_{>1},
\end{equation*}
also streng monoton fallend. Grenzwerte: für $x\nearrow1$ geht $x^2-1\to
0^+$ mit Zähler $2x\to2>0$, also $f\to+\infty$; für $x\nearrow-1$ geht
$x^2-1\to0^+$, Zähler $2x\to-2<0$, also $f\to-\infty$. Und $f(0)=0$.

\begin{ergebnis}
\begin{equation*}
  \text{a) } f(x)=1+9\mathrm e^{-(x-2)^2};\quad
  \text{b) } g(x)=-\ln(x-1)+(\ln2-1)x+2-2\ln2;\quad
  \text{c) } f(x)=\frac{2x}{x^2-1}.
\end{equation*}
\end{ergebnis}

\begin{kontrolle}
\textbf{Probe (a):} $1+9\mathrm e^0=10$; beide Grenzwerte $1$ wegen
$\mathrm e^{-(x-2)^2}\to0$. \checkmark

\textbf{Probe (b):} $g(2)=2(\ln2-1)+2-2\ln2=0$;
$g(3)=3(\ln2-1)-\ln2+2-2\ln2=-1$. Die Pole: $-\ln(x-1)\to+\infty$ bei
$x\searrow1$; der Grenzwert $-\infty$ kommt von $-(\ln2-1)\,x\to-\infty$
plus $-\ln(x-1)\to-\infty$. \checkmark

\textbf{Probe (c):} $f(0)=0$; $f'(x)<0$ auf $R_{>1}$; Pole wie gefordert.
Zahlencheck: für $x=2$ ist $f(2)=4/3>0$ und für $x\to1^+$ unendlich
positiv, für grosse $x$ geht $f\to0^+$ -- eine stetig fallende, positive
Kurve ohne Extremum. \checkmark
\end{kontrolle}

\begin{merksatz}
Eine elementare Funktion baut man aus Grundbausteinen (Pol, Glocke,
$\ln$, lineare Anpassung), die man additiv kombiniert, um mehrere
Bedingungen gleichzeitig zu erfüllen. Jede Bedingung einzeln nachrechnen;
Dynamik-Baustein (Pol) und Wert-Anpassung (lineare Funktion) getrennt
halten.
\end{merksatz}
